\documentclass[11pt]{article}
\usepackage[a4paper,margin=27mm]{geometry}
\usepackage[T1]{fontenc}
\usepackage{lmodern,microtype,amsmath,amssymb,amsthm,mathtools,enumitem,booktabs}
\usepackage[hidelinks]{hyperref}
\newtheorem{theorem}{Theorem}
\newtheorem{lemma}[theorem]{Lemma}
\newtheorem{proposition}[theorem]{Proposition}
\newtheorem{corollary}[theorem]{Corollary}
\theoremstyle{remark}\newtheorem{remark}[theorem]{Remark}
\newtheorem{example}[theorem]{Example}
\newcommand{\Dscr}{\mathcal D}
\newcommand{\M}{\mathcal M}
\newcommand{\R}{\mathbb R}
\newcommand{\K}{\mathcal K}
\newcommand{\supp}{\operatorname{supp}}
\newcommand{\conv}{\operatorname{conv}}
\newcommand{\cl}{\operatorname{cl}}
\newcommand{\one}{\mathbf 1}
\title{Conditional Shunning of Error}
\author{}\date{12 October 2026}
\begin{document}\maketitle

\noindent
Fix a set $T$ of worlds. The systems $\Sigma_T$ below model an evaluator
with conditional values: conditional on $T$ she only shuns error, and
conditional on the complement of $T$ she also seeks truth, so that (G) is
strict there. Unanimous problems are scored strictly at every world. For
$T$ equal to all worlds, $\Sigma_T$ is the accepted exceptional system for
the vacuous model $V$; for $T$ empty it is the strict two-slope system. We show that every $\Sigma_T$
satisfies (IB), (LA), (C), (S), and the guidance order weakened exactly at
the worlds of $T$. Closed competitive selection holds for every $T$, but
only because the selected model need not be Bayes at its own prior: no
selection of Bayes models works once $n\ge3$ and $\varnothing\ne T\ne\Omega$.
We identify the undominated models of each system and the credal sets
covered by the whole family. They are the normalized boxes in which some
lower bounds may vanish, and an upper bound may be infinite where the lower
bound vanishes. The family covers $V$, every face, every precise model,
every credal set in two worlds, and many proper targets containing three
vertices. No system in the family covers $\varepsilon$-contamination
classes, or any polytope with a facet not of the form $\{m_k=0\}$ or
$\{m_b=\rho m_a\}$. Every result below is proved; the open questions are
listed at the end. The rankings are complete. The note was checked by an
independent referee, and the referee's corrections are incorporated; no
human verification is asserted.

\section{The systems}

There are $n\ge2$ worlds, $\Omega=\{1,\ldots,n\}$, with simplex $\Delta$,
vertices $e_i$ and faces $\Delta_S=\{m\in\Delta:\supp m\subseteq S\}$.
A model is a measurable acceptance set, identified modulo Gaussian null
sets. Each binary problem $\{g,-g\}$ is counted once by $g\in H=\{g_1>0\}$,
which carries the normalized Gaussian measure $\mu$. It is convenient to
write $\bar\mu$ for the centred Gaussian measure on $\R^n$ scaled so that
$\bar\mu(H)=1$: a set $A$ of orientations with $A\cap(-A)=\varnothing$
represents a set of problems of $\mu$-measure $\bar\mu(A)$. Coordinate
hyperplanes are null and are ignored. A problem is \emph{unanimous} if $g$
or $-g$ is strictly positive, and \emph{mixed} otherwise. The advice of a
model on $g\in H$ is $+$ (accept $g$ only), $0$ (silence), $-$ (accept $-g$
only) or $\times$ (accept both). At world $i$ it is good, silent, bad or
contradictory. We write $x^+=\max(x,0)$ and $x^-=\max(-x,0)$, and use
$\M_K=\{g:\min_{m\in K}m\cdot g>0\}\cup\{g\ge0:g\ne0\}$ for a credal set
$K$, so that $V=\M_\Delta$.

Fix $T\subseteq\Omega$ and put $U=\Omega\setminus T$. Choose numbers
\[
 0<\lambda_j\le\eta_j<\infty\ (j\in U),\qquad
 \lambda_i\coloneqq0<\eta_i<\infty\ (i\in T).
\]
At world $i$ the local penalties are
\[
\begin{array}{c|cccc}
 &\text{good}&\text{silence}&\text{bad}&\text{contradictory}\\\hline
 \text{mixed}&0&\lambda_i|g_i|&(\lambda_i+\eta_i)|g_i|
     &2(\lambda_i+\eta_i)|g_i|\\
 \text{unanimous}&0&|g_i|&2|g_i|&4|g_i|.
\end{array}
\tag{1}
\]
The loss $L_i(X)$ is the $\mu$-integral of the penalty of the advice of $X$.
Smaller is better, and $\Sigma_T=\Sigma_T(\lambda,\eta)$ denotes the
resulting complete rankings. For $T=\Omega$ the mixed row is
$0,0,\eta_i|g_i|,2\eta_i|g_i|$, which is the accepted exception. For
$T=\varnothing$ it is the strict table (3) of \emph{Coverage after Neutral
Abstention}.

\begin{proposition}\label{prop:basic}
Each $\Sigma_T$ has finite losses on all models and satisfies (IB). It
satisfies (G) in the following weakened form. Every pointwise improvement
of advice weakly lowers the loss. It lowers the loss strictly unless, up
to a null set, every strict improvement is a change from silence to good
advice on a mixed problem at a world of $T$. Such changes leave the loss
unchanged.
\end{proposition}
\begin{proof}
All penalties are at most $4\max(1,\lambda_i+\eta_i)|g_i|$, which is
integrable. Shared advice contributes the same amount to both losses, which
gives (IB). In (1), every adjacent gap is positive except good versus
silence on mixed problems at $i\in T$, where both penalties are $0$.
\end{proof}

\section{Expected penalties}

For a prior $r\in\Delta$ and a mixed $g\in H$, write $E(a)$ for the
$r$-expected penalty of advice $a$, and put
\[
 \psi_r(g)=\sum_ir_i\bigl(\lambda_ig_i^+-\eta_ig_i^-\bigr)
 =\sum_ir_i\min(\lambda_ig_i,\eta_ig_i).
\]
A direct computation from (1) gives
\[
 E(0)-E(+)=\psi_r(g),\qquad E(0)-E(-)=\psi_r(-g),\qquad
 E(\times)=2\bigl(E(+)+E(-)\bigr).
\tag{2}
\]
Here $E(+)+E(-)=\sum_ir_i(\lambda_i+\eta_i)|g_i|>0$, so contradiction is
never optimal, and $\psi_r(g)+\psi_r(-g)=-\sum_ir_i(\eta_i-\lambda_i)|g_i|
\le0$. On a unanimous problem good advice is the unique optimum.

It is useful to separate the two kinds of world. For
$\alpha\in[0,\infty]^T$ and $\beta\in[0,\infty)^U$ with $\beta\ne0$, put
\[
 Q_{\alpha,\beta}(g)=\sum_{j\in U}\beta_j\bigl(\lambda_jg_j^+-\eta_jg_j^-\bigr)
 -\sum_{i\in T}\alpha_i\eta_ig_i^-\in[-\infty,\infty),
 \qquad \infty\cdot0\coloneqq0,
\]
\[
 M(\alpha,\beta)=\{g:Q_{\alpha,\beta}(g)>0\}\cup\{g\ge0:g\ne0\},
 \qquad I(\alpha)=\{i\in T:\alpha_i=\infty\}.
\]
Thus $M(\alpha,\beta)$ rejects every $g$ with $g_i<0$ for some
$i\in I(\alpha)$. Both $Q$ and $M$ are unchanged when $(\alpha,\beta)$ is
multiplied by a positive number, so we normalize $\beta\in\Delta_U$
whenever $U\ne\varnothing$. If $s=r(U)>0$, then
$\psi_r=s\,Q_{r_T/s,\,r_U/s}$.

\begin{lemma}\label{lem:Q}
Let $\beta\ne0$.
\begin{enumerate}[label=(\alph*),leftmargin=*]
\item $Q_{\alpha,\beta}$ is concave, positively homogeneous and
nondecreasing in each coordinate. Hence $M(\alpha,\beta)$ is coherent, and
$Q_{\alpha,\beta}(g)+Q_{\alpha,\beta}(-g)\le0$.
\item The set $\{Q_{\alpha,\beta}=0\}$ is null.
\item The advice of $M(\alpha,\beta)$ is never contradictory and is good on
every unanimous problem.
\item The map $(\alpha,\beta)\mapsto L(M(\alpha,\beta))$ is continuous on
$[0,\infty]^T\times\Delta_U$, where $[0,\infty]$ has its usual compact
topology.
\end{enumerate}
\end{lemma}
\begin{proof}
(a) Each term $\beta_j\min(\lambda_jg_j,\eta_jg_j)$ is a minimum of linear
functions because $\lambda_j\le\eta_j$. Each term $-\alpha_i\eta_ig_i^-$ is
concave, including when $\alpha_i=\infty$, where it is $0$ for $g_i\ge0$ and
$-\infty$ otherwise. All terms are positively homogeneous and
nondecreasing. Concavity and homogeneity give
$Q(g+h)\ge Q(g)+Q(h)$. So $\{Q>0\}$ is a convex cone, it omits $0$, and
adding a nonnegative gamble cannot leave it. Hence $M(\alpha,\beta)$ is
coherent, and $Q(g)+Q(-g)\le Q(0)=0$.

(b) On an open orthant, $Q$ is identically $-\infty$ if some $i\in
I(\alpha)$ has $g_i<0$. Otherwise it is linear, and its coefficient on
$g_j$, for $j\in\supp\beta$, is $\beta_j\lambda_j$ or $\beta_j\eta_j$, both
positive. Its zero set in the orthant is therefore contained in a
hyperplane.

(c) Both $g$ and $-g$ lie in $\{Q>0\}$ for no $g$, by (a); and if $g\ge0$ then $Q(-g)\le Q(0)=0$ by monotonicity. If $g\gg0$ then $g\in M$, and
$Q(-g)=-\sum_U\beta_j\eta_jg_j-\sum_T\alpha_i\eta_ig_i<0$, so $-g\notin M$.

(d) Fix $g$ with no zero coordinate. Each term of $Q_{\alpha,\beta}(g)$ is
continuous in $(\alpha,\beta)$ with values in $[-\infty,\infty)$: terms with
$g_i>0$ do not depend on $\alpha_i$, and $\alpha_i\mapsto\alpha_i\eta_i|g_i|$
is continuous into $[0,\infty]$. If $(\alpha^k,\beta^k)\to(\alpha,\beta)$
and $Q_{\alpha,\beta}(\pm g)\ne0$, the signs of $Q_{\alpha^k,\beta^k}(\pm g)$
are eventually those of $Q_{\alpha,\beta}(\pm g)$. By (b) this holds for
almost every $g$. The advice therefore converges almost everywhere, and
dominated convergence with the bound in Proposition~\ref{prop:basic}
gives convergence of the losses.
\end{proof}

\begin{proposition}\label{prop:bayes}
Let $r\in\Delta$.
\begin{enumerate}[label=(\alph*),leftmargin=*]
\item If $r\notin\Delta_T$, the $r$-Bayes model is unique modulo null sets
and equals $M(r_T/s,r_U/s)$, where $s=r(U)$. It is coherent and
undominated. Its credal set is the normalized box
\[
 K_r=\Bigl\{\frac{a}{\one\cdot a}:
 r_i\lambda_i\le a_i\le r_i\eta_i\text{ for all }i\Bigr\},
\]
whose lower bounds vanish on $T\cup(\Omega\setminus\supp r)$.
\item If $r\in\Delta_T$ and $S=\supp r$, then $\min_X r\cdot L(X)=0$, and
$X$ is $r$-Bayes exactly when $L_i(X)=0$ for all $i\in S$. Equivalently,
for almost every problem, $X$ gives good advice if it is unanimous, is
silent if both orientations have coordinates of both signs on $S$, and
otherwise either accepts the orientation $h$ with $h_S\gg0$ or is silent.
If $S\ne\Omega$, these choices are non-null. In particular
$\{\psi_r>0\}\cup\{g\ge0:g\ne0\}=V$, whose credal set is $\Delta$ and not
the normalized box $\Delta_S$.
\end{enumerate}
\end{proposition}
\begin{proof}
(a) Some $j\in U$ has $r_j>0$, so by Lemma~\ref{lem:Q}(b) the tie set
$\{\psi_r=0\}$ is null. By (2), the pointwise optimal advice is unique
almost everywhere, and it is the advice of $M(r_T/s,r_U/s)$ by
Lemma~\ref{lem:Q}(c). Integrating the nonnegative difference from the
pointwise optimum shows that a model is $r$-Bayes exactly when it agrees
with this advice almost everywhere. A weak dominator of a model that is
uniquely $r$-Bayes is itself $r$-Bayes, so it agrees with that model and
ties with it. Hence the model is undominated. For the credal set, note that
$\psi_r(g)=\min\{a\cdot g:r_i\lambda_i\le a_i\le r_i\eta_i\}$. Every point
$a$ of this box is nonzero, since $a_j\ge r_j\lambda_j>0$, and
$a\cdot g=(\one\cdot a)\,(a/\one\cdot a)\cdot g$. So $\psi_r(g)>0$ exactly
when $m\cdot g>0$ for all $m\in K_r$. The set $K_r$ is compact and convex,
since the normalization map sends segments to segments.

(b) The model $V$ has $L_i(V)=0$ for every $i\in T$. Hence
$\min_Xr\cdot L(X)=0$, and $r$-Bayes means $L_i(X)=0$ for all $i\in S$. At
$i\in T$ a mixed problem has positive penalty exactly for bad or
contradictory advice, and a unanimous problem has positive penalty for all
advice except good advice. Combining the conditions over $i\in S$ gives
the description. If $S\ne\Omega$, the problems with an orientation $h$
satisfying $h_S\gg0$ and $h_k<0$ for some $k\notin S$ form an open set. As
$\psi_r=-\sum_S r_i\eta_ig_i^-\le0$, the strict set is $V$.
\end{proof}

\begin{remark}\label{rem:variant}
One might let the worlds of $T$ also credit gambles that pay throughout
$T$. In this variant, fix $\kappa_i>0$ for $i\in T$. On a mixed problem
whose $T$-coordinates all have the same sign ($g_T\gg0$ or $g_T\ll0$), world
$i\in T$ uses the row $0,\ \kappa_i|g_i|,\ (\kappa_i+\eta_i)|g_i|,\
2(\kappa_i+\eta_i)|g_i|$. On other mixed problems it uses
$0,\ 0,\ \eta_i|g_i|,\ 2\eta_i|g_i|$. Worlds outside $T$ and unanimous
problems are scored as in (1). This variant can be incoherent. Take
$n=3$, $T=\{1,2\}$, $\kappa_i=\eta_i=\lambda_3=\eta_3=1$ and the uniform
prior. The expected costs of $(+,0,-)$, multiplied by $3$, are
$(6,7,8)$ at $(2,2,-3)$, $(10,10.5,21.01)$ at $(-10,0.01,10.5)$, and
$(4,3.75,8.505)$ at their midpoint $(-4,1.005,3.75)$. On each open
orthant all differences of expected costs are linear forms, with a
nonzero coefficient on $g_3$, so ties are null. The uniform-prior Bayes
model is therefore unique modulo null sets. By continuity it accepts open
neighbourhoods $N_1,N_2$ of the first two gambles, and is silent on an
open neighbourhood of their midpoint. No coherent model agrees with it
almost everywhere. If $D$ did, $D$ would contain full-measure subsets
$N_1',N_2'$, and every point $x$ of $\tfrac12(N_1+N_2)$ is a midpoint of
some $a\in N_1'$ and $b\in N_2'$, since $\{a\in N_1:2x-a\in N_2\}$ is a
nonempty open set. So $D$ would contain an open neighbourhood of the
midpoint. On the other hand, (S) yields a coherent Bayes model at every
prior, by the simplex argument of part~(a) of the first lemma in
Appendix~C of \emph{Full Coverage by Homothetic Credal Families v3}. Hence
the variant violates (S).
\end{remark}

\section{The undominated models}

For $S\subseteq T$, let $R_S$ be the open cone of mixed $h$ with $h_S\gg0$.
For $k\in\Omega$ put
\[
 c_k(h)=-\lambda_kh_k^++\eta_kh_k^-.
\]
This is the change in penalty at world $k$ when a model that was silent on
the problem of $h$ accepts $h$.

\begin{theorem}\label{thm:undominated}
If $U\ne\varnothing$, a model is undominated in $\Sigma_T$ exactly when it
agrees almost everywhere with $M(\alpha,\beta)$ for some
$\alpha\in[0,\infty]^T$ and $\beta\in\Delta_U$. If $U=\varnothing$, the only
undominated model is $V$. In either case every undominated model is
coherent modulo null sets, so $\Sigma_T$ satisfies (LA) and (C).
\end{theorem}
\begin{proof}
\emph{Sufficiency.} If $I=I(\alpha)$ is empty, $M(\alpha,\beta)$ is the
unique Bayes model at the prior proportional to $(\alpha,\beta)$, and
Proposition~\ref{prop:bayes}(a) applies. Suppose $I\ne\varnothing$. By
Lemma~\ref{lem:Q}, $M(\alpha,\beta)$ accepts a mixed $g$ only if $g_I\ge0$,
so $L_i(M(\alpha,\beta))=0$ for $i\in I$. Let $Y$ weakly dominate it. Then
$L_I(Y)=0$, and Proposition~\ref{prop:bayes}(b) for a prior with support
$I$ fixes the advice of $Y$ except on $R_I$, where $Y$ accepts some set
$A_Y$. Put $w=(\alpha_{T\setminus I},\beta)$, a finite nonzero weight on
$\Omega\setminus I$. On $R_I$ we have
$-\sum_{k\notin I}w_kc_k=Q_{\alpha,\beta}$, so every model $Z$ with
$L_I(Z)=0$ satisfies
\[
 \sum_{k\notin I}w_kL_k(Z)=\sum_{k\notin I}w_kL_k(V)
 -\int_{A_Z}Q_{\alpha,\beta}\,d\bar\mu.
\]
By Lemma~\ref{lem:Q}(b) this is minimized exactly when
$A_Z=R_I\cap\{Q_{\alpha,\beta}>0\}$ modulo null sets, which is the choice
made by $M(\alpha,\beta)$. Weak dominance gives
$\sum w_kL_k(Y)\le\sum w_kL_k(M(\alpha,\beta))$, so $Y$ is also a minimizer
and agrees with $M(\alpha,\beta)$ almost everywhere. Hence nothing
dominates $M(\alpha,\beta)$.

\emph{Necessity.} Let $X$ be undominated. The closure of the set of
attainable loss vectors is convex (Gaussian measure has no atoms; see
part~(b) of the first lemma in Appendix~C of \emph{Full Coverage by
Homothetic Credal Families v3}). Separating it from $L(X)-(0,\infty)^n$ shows that $X$ is
$r$-Bayes for some $r\in\Delta$. If $r\notin\Delta_T$,
Proposition~\ref{prop:bayes}(a) identifies $X$. Otherwise the set
$S=\{i\in T:L_i(X)=0\}$ contains $\supp r$ and is nonempty. By
Proposition~\ref{prop:bayes}(b), $X$ is fixed except on $R_S$, where it
accepts a set $A_X$. For $k\notin S$,
\[
 L_k(Z)=L_k(V)+\int_{A_Z}c_k\,d\bar\mu
\]
for every model $Z$ with $L_S(Z)=0$. Thus $A_X$ is Pareto optimal for the
vector $(\int_Ac_k\,d\bar\mu)_{k\notin S}$ over measurable $A\subseteq R_S$,
or $X$ would be dominated. By Lyapunov's theorem these vectors form a
compact convex set. Separation gives $w\in\Delta_{\Omega\setminus S}$ such
that $A_X$ minimizes $\int_A\sum_kw_kc_k\,d\bar\mu$. Note that
$-\sum_kw_kc_k=\psi_w$. If $\supp w$ meets $U$, then by
Lemma~\ref{lem:Q}(b) $A_X=R_S\cap\{\psi_w>0\}$ modulo null sets. Hence
$X=M(\alpha,\beta)$ almost everywhere, where $\alpha_i=\infty$ on $S$,
$\alpha_i=w_i/w(U)$ on $T\setminus S$, and $\beta=w_U/w(U)$. If
$\supp w\subseteq T\setminus S$, then $\psi_w(h)<0$ unless
$h_{\supp w}\ge0$. So almost every accepted $h\in R_S$ has
$h_{\supp w}\ge0$, and $L_k(X)=0$ for $k\in\supp w$. This contradicts the
definition of $S$. This argument needs $S\ne\Omega$; if $U\ne\varnothing$
this holds because $S\subseteq T$. If $U=\varnothing$ and $S\ne\Omega$, then
$\supp w\subseteq T\setminus S$, and the second alternative gives a
contradiction. Hence $S=\Omega$, $R_\Omega=\varnothing$, and $X=V$.

Every $M(\alpha,\beta)$ is coherent by Lemma~\ref{lem:Q}(a), giving (LA).
A model tied at every world with an undominated model is undominated,
hence coherent modulo null sets. This gives (C).
\end{proof}

\begin{proposition}\label{prop:V}
$V$ is undominated in $\Sigma_T$ if and only if $|U|\le1$.
\end{proposition}
\begin{proof}
If $U=\varnothing$, Theorem~\ref{thm:undominated} applies. If $U=\{j\}$, then
$M(\alpha,e_j)$ with $\alpha=\infty$ on $T$ accepts exactly the gambles
with $g_T\ge0$ and $g_j>0$, up to null sets, which is $V$.

Let $j\ne k$ lie in $U$. For $0<\delta<1$ let $A$ be the open cone of $g$
with $g_j>0$, $-\delta g_j<g_k<-\tfrac\delta2g_j$, and
$0<g_l<\delta^2g_j$ for $l\notin\{j,k\}$. Let $B$ be its image under the
transposition of $j$ and $k$. Put $Y=V\cup A\cup B$. On $A\cup B$ the model
$V$ is silent and $Y$ accepts. At worlds of $T$ this changes silence to good
advice on mixed problems, which costs nothing. At $l\in U\setminus\{j,k\}$
it is a strict gain. At world $j$ the change is
\[
 -\lambda_j\int_Ag_j\,d\bar\mu+\eta_j\int_B|g_j|\,d\bar\mu
 \le(\eta_j\delta-\lambda_j)\int_Ag_j\,d\bar\mu,
\]
because the transposition preserves $\bar\mu$ and $|g_k|\le\delta g_j$ on
$A$. For $\delta<\min(\lambda_j/\eta_j,\lambda_k/\eta_k)$ this is negative,
and world $k$ is symmetric. The sets $A$, $B$, $-A$, $-B$ and $\pm V$ meet
only in null sets, so $Y$ gives no contradictory advice. Hence $Y$
dominates $V$.
\end{proof}

\section{Closed competitive selection}

A model that is $r$-Bayes is competitive on every $S\supseteq\supp r$:
against any $X$, the inequality $r\cdot L(C)\le r\cdot L(X)$ needs some
$i\in\supp r$ with $L_i(C)\le L_i(X)$. Conversely, separation in the
coordinates of $S$ shows that a model competitive on $S$ is Bayes for some
prior supported in $S$. So (S) asks for coherent undominated models $C_q$
that are Bayes at \emph{some} prior supported in $\supp q$, not necessarily
at $q$.

\begin{remark}\label{rem:lsc}
For a fixed world $i$, the values $L_i(X)$ fill the interval from $0$ to
$L_i$ of the everywhere contradictory model. To see this, interpolate
between $D_i=\{g_i>0\}$ and that model on growing balls. Hence the sections
$\{q:C_q\preceq_iX\}$ are closed for all $X$ if and only if
$q\mapsto L_i(C_q)$ is lower semicontinuous.
\end{remark}

\begin{theorem}\label{thm:S}
Every $\Sigma_T$ satisfies (S) with a selection whose losses are
continuous in the prior.
\end{theorem}
\begin{proof}
If $U=\varnothing$, take $C_q=V$. If $T=\varnothing$, take the Bayes
selection of Proposition~\ref{prop:bayes}(a). Its losses are continuous by
Lemma~\ref{lem:Q}(d). Otherwise let $t=|T|$, fix $0<\varepsilon\le\tfrac12$
and $d\in\Delta_U$. For $q\in\Delta$ put $s=q(U)$ and
\[
 m(s)=\max(s,\varepsilon),\qquad
 k(s)=t\min\bigl(1,\max(0,2-s/\varepsilon)\bigr),\qquad
 w(s)=\max(0,1-s/\varepsilon).
\]
Define
\[
 \alpha_i(q)=
 \begin{cases}
 q_i/(m(s)-k(s)q_i),& k(s)q_i<m(s),\\
 \infty,&\text{otherwise},
 \end{cases}
 \quad(i\in T),\qquad
 \beta(q)=\frac{q_U+w(s)\,d}{s+w(s)},
\]
and let $C_q=M(\alpha(q),\beta(q))$.

\emph{Continuity.} Since $m\ge\varepsilon>0$, the denominator
$m-kq_i$ tends to $0$ only when $q_i\ge\varepsilon/t$. Hence
$q\mapsto\alpha(q)\in[0,\infty]^T$ is continuous, and $\alpha_i(q)=0$
exactly when $q_i=0$. The denominator $s+w(s)$ is positive, so
$\beta(q)\in\Delta_U$ is continuous. Lemma~\ref{lem:Q}(d) makes every
$L_i(C_q)$ continuous, and Remark~\ref{rem:lsc} gives closed sections.

\emph{Competitiveness.} If $I(\alpha(q))\ne\varnothing$, each of its
indices has $q_i>0$. Then $C_q$ has zero loss at that world, as in the
proof of Theorem~\ref{thm:undominated}, so it is Bayes at a vertex in
$\Delta_{\supp q}$. If $I(\alpha(q))=\varnothing$, then $s>\varepsilon$:
for $s\le\varepsilon$ we have $k=t$ and $m=\varepsilon$, and some $i\in T$
has $q_i\ge(1-\varepsilon)/t$, so $kq_i\ge1-\varepsilon\ge m$. Hence $w=0$,
$\beta=q_U/s$, and $C_q$ is the unique Bayes model at the prior
proportional to $(\alpha(q),q_U/s)$, whose support is $\supp q$.

By Theorem~\ref{thm:undominated} every $C_q$ is coherent and undominated.
For $s\ge2\varepsilon$ the selection is the Bayes selection.
\end{proof}

The freedom to select a non-Bayes model is essential.

\begin{proposition}\label{prop:bayessel}
A selection with $C_q$ being $q$-Bayes for every $q$ satisfies (S) if and
only if $T=\varnothing$, $T=\Omega$, or $n=2$.
\end{proposition}
\begin{proof}
The three cases are covered by the selections above. When $n=2$ and
$T=\{i\}$, take $C_{e_i}=V$; it is the limit of the Bayes models from the
only direction of approach. Otherwise $U\ne\varnothing$, and some $p\in\Delta_T$ has
$S=\supp p\ne\Omega$ and admits two directions of approach: if $|T|\ge2$,
take $p=e_i$ with $i\in T$, so $T\setminus S\ne\varnothing$; if $|T|=1$,
then $|U|\ge2$. For $a\in[0,\infty)^{T\setminus S}$ and $b\in\Delta_U$, the
priors $q_\tau\propto(p_S/\tau,a,b)$ lie outside $\Delta_T$. By
Lemma~\ref{lem:Q}(d) their Bayes models $M((p_S/\tau,a),b)$ converge in
loss, as $\tau\to0$, to $M_{a,b}\coloneqq M(\alpha^*,b)$ with
$\alpha^*=\infty$ on $S$ and $\alpha^*=a$ on $T\setminus S$, while
$q_\tau\to p$. Lower semicontinuity requires $L(C_p)\le L(M_{a,b})$. Since
$M_{a,b}$ is undominated, the two loss vectors are equal. So all $M_{a,b}$
have the same losses.

Now take $(a,b)\ne(a',b')$. If $b\ne b'$, choose $j,k\in U$ with
$b_j>b'_j$ and $b_k<b'_k$. On gambles with $g_S\gg0$, $g_j>0>g_k$ and
the other $U$-coordinates positive and sufficiently small, and the
remaining coordinates positive, $Q$ for $(a,b)$ is positive and $Q$ for
$(a',b')$ is negative whenever $g_j/|g_k|$ lies strictly between
$b_k\eta_k/(b_j\lambda_j)$ and $b'_k\eta_k/(b'_j\lambda_j)$ (the latter is
$\infty$ if $b'_j=0$). If $b=b'$ and $a_i<a'_i$, compare $g_j>0$
with $g_i<0$ for some $j\in\supp b$ in the same way. In each case the two
models differ on an open subset of $R_S$. By the sufficiency part of
Theorem~\ref{thm:undominated}, with $I=S$, each is the unique minimizer
modulo null sets of a weighted loss among models with $L_S=0$. If the two
had equal losses, each would minimize the other's weighted loss, which is
impossible. This contradiction proves the claim.
\end{proof}

\section{The covered credal sets}

Call $(l,u)\in[0,\infty)^n\times[0,\infty]^n$ \emph{admissible} if
$l_i\le u_i$ for every $i$, $u_i<\infty$ whenever $l_i>0$, and $u\ne0$.
Put $P=\{i:u_i>0\}$, $L=\{i:l_i>0\}$ and $F=\{i:u_i<\infty\}$. Define
\[
 K(l,u)=\cl\Bigl\{\frac{a}{\one\cdot a}:a\in\R^n,\ a\ne0,\
 l_i\le a_i\le u_i\text{ for all }i\Bigr\}.
\]
Thus each world is either \emph{anchored} ($0<l_i\le u_i<\infty$) or
\emph{free} ($l_i=0$), and a free world may have an infinite upper bound.
When every world in $P$ is anchored, these are the normalized boxes of
\emph{Coverage after Neutral Abstention}, equation (4).

\begin{lemma}\label{lem:intrinsic}
For admissible $(l,u)$,
\[
 K(l,u)=\{m\in\Delta_P:\ l_am_b\le u_bm_a\ \text{ for all }a\in L,\
 b\in P\cap F\}.
\]
That is, $m_b/m_a\le u_b/l_a$ for all anchored $a$ and all $b$ with
finite upper bound. If $L=\varnothing$, then $K(l,u)=\Delta_P$. Moreover
$e_k\in K(l,u)$ if and only if $k\in P$ and either $u_k=\infty$ or
$L\subseteq\{k\}$.
\end{lemma}
\begin{proof}
Write $K'$ for the right side. If $l\le a\le u$ with $a$ finite, then
$l_aa_b\le a_au_b$, so normalized box points lie in the closed set $K'$.
For the converse, let $m\in K'$. If $L=\varnothing$, $sm$ lies in the box
for small $s>0$. If $m_a>0$ for all $a\in L$, put $s=\max_{a\in L}l_a/m_a$.
Then $sm\ge l$, and for $b\in F$ some $a\in L$ gives
$sm_b=(l_a/m_a)m_b\le u_b$. So $m$ is a normalized box point. If $m_a=0$
for some $a\in L$, the constraints force $m_b=0$ for every $b\in P\cap F$,
so $\supp m\subseteq\{u=\infty\}$. Fix a box point $a^0$; then $a^0+Nm$
lies in the box for $N>0$, and its normalization tends to $m$. For the
vertex criterion, $m=e_k$ violates the constraint for $(a,b)=(a,k)$
exactly when $a\in L\setminus\{k\}$ and $k\in F$.
\end{proof}

\begin{lemma}\label{lem:credal}
For $\alpha\in[0,\infty]^T$ and $\beta\in\Delta_U$, $M(\alpha,\beta)$
agrees almost everywhere with $\M_K$, where $K=K(l,u)$, $l_j=\beta_j\lambda_j$
and $u_j=\beta_j\eta_j$ on $U$, and $l_i=0$ and $u_i=\alpha_i\eta_i$ on $T$.
The two sets are equal when $I(\alpha)=\varnothing$.
\end{lemma}
\begin{proof}
Coordinatewise, $Q_{\alpha,\beta}(g)=\inf\{a\cdot g:a\text{ finite},\
l\le a\le u\}$, since $\inf_{0\le a_i<\infty}a_ig_i$ is $0$ or $-\infty$.
Let $I=I(\alpha)$ and let $K'$ be the normalized compact box with
$a_I=0$, as in the proof of Proposition~\ref{prop:bayes}(a). Then $Q>0$
means $g_I\ge0$ together with $\min_{K'}m\cdot g>0$. Normalizing
$a'+a_I$ produces convex combinations of points of $K'$ and $\Delta_I$,
with any weight in $(0,1]$ on $K'$; hence $K=\conv(K'\cup\Delta_I)$. So
$\min_Km\cdot g>0$ means $g_I\gg0$ together with $\min_{K'}m\cdot g>0$.
The two sets differ only where some $g_i=0$ with $i\in I$.
\end{proof}

\begin{lemma}\label{lem:ident}
Let $K,K'\subseteq\Delta$ be nonempty, compact and convex. If $\M_K$ and
$\M_{K'}$ agree almost everywhere, then $K=K'$.
\end{lemma}
\begin{proof}
Suppose some $m^0\in K'\setminus K$. Separating $m^0$ from $K$ within the
affine hull of $\Delta$ gives $c$ and $d$ with $c\cdot m>d>c\cdot m^0$ for
all $m\in K$. Put $g=c-d\one$. Then $\min_Km\cdot g>0>\min_{K'}m\cdot g$.
Both minima are continuous in $g$, so these strict inequalities hold on an
open neighbourhood of $g$. No point $h$ of that neighbourhood is
nonnegative, since $m^0\cdot h<0$. The neighbourhood therefore lies in
$\M_K\setminus\M_{K'}$, which is impossible. By symmetry, $K\subseteq K'$
as well.
\end{proof}

\begin{theorem}\label{thm:coverage}
Let $\K$ be the set of all $K(l,u)$ with $(l,u)$ admissible. Then:
\begin{enumerate}[label=(\alph*),leftmargin=*]
\item Fix $T$ and $(\lambda,\eta)$. If $U\ne\varnothing$, the coherent
undominated models of $\Sigma_T$ are, modulo null sets, exactly the
$\M_{K(l,u)}$ with $(l,u)$ admissible, $l=0$ on $T$, $u_T\in[0,\infty]^T$,
and $(l_j,u_j)=c_j(\lambda_j,\eta_j)$ on $U$ for some $c\in[0,\infty)^U$
with $c\ne0$. If $U=\varnothing$, the only one is $V$.
\item The family of all $\Sigma_T(\lambda,\eta)$, every one of which
satisfies the conclusions of Proposition~\ref{prop:basic},
Theorem~\ref{thm:undominated} and Theorem~\ref{thm:S}, covers exactly
the targets $\M_K$ with $K\in\K$.
\item Suppose $K=K(l,u)$ with $L\ne\varnothing$.
\begin{enumerate}[label=(c\arabic*),leftmargin=*]
\item If $u$ is finite and $p$ is any prior with $\supp p=P$, then $\M_K$
is the unique Bayes model at $p$ in $\Sigma_T$ with $T=\Omega\setminus L$,
$\lambda_j=l_j/p_j$ and $\eta_j=u_j/p_j$ on $L$, and $\eta_i=u_i/p_i$ on
$P\setminus L$; any $\eta_i>0$ may be used off $P$.
\item If $I\coloneqq\{i:u_i=\infty\}$ is nonempty, then in the system used
in the proof of (b), $\M_K$ agrees almost everywhere with $M(\alpha,\beta)$,
where $I(\alpha)=I$, and it is Bayes exactly at the priors in $\Delta_I$.
\end{enumerate}
\end{enumerate}
\end{theorem}
\begin{proof}
(a) By Theorem~\ref{thm:undominated} and Lemma~\ref{lem:credal}, the
undominated models are, modulo null sets, the $\M_{K(l,u)}$ with
$l_j=\beta_j\lambda_j$ and $u_j=\beta_j\eta_j$ on $U$, and $l_i=0$ and
$u_i=\alpha_i\eta_i$ on $T$. Since $K(l,u)$ is unchanged when $(l,u)$ is
multiplied by a positive number, any $c\ne0$ arises, with $\beta=c/\sum_Uc$
and $\alpha_i=u_i/(\eta_i\sum_Uc)$. For (b),
given admissible $(l,u)$ with $L\ne\varnothing$, put $T=\Omega\setminus L$,
$\beta=l_L/\sum_Ll$, $\lambda_j=\sum_Ll$ and $\eta_j=(u_j/l_j)\sum_Ll$ on
$L$, and $\eta_i=1$, $\alpha_i=u_i$ on $T$. Lemma~\ref{lem:credal} returns
$K(l,u)$. If $L=\varnothing$, then $K(l,u)=\Delta_P$, which is also
$K(l',u')$ with one anchored $j\in P$ ($l'_j=u'_j=1$), $u'=\infty$ on
$P\setminus\{j\}$ and $u'=0$ off $P$. Conversely, if $\M_K$ is
undominated in some $\Sigma_T$, then by (a) it agrees almost everywhere
with some $\M_{K(l,u)}$, or with $V=\M_\Delta$, and Lemma~\ref{lem:ident}
gives $K=K(l,u)$ or $K=\Delta$. Both lie in $\K$.

(c1) With these parameters the box of Proposition~\ref{prop:bayes}(a) at
$p$ is $\prod_i[l_i,u_i]$. Here $L\subseteq\supp p$ meets $U$, so
$p\notin\Delta_T$, and the Bayes model is unique.

(c2) The model $M(\alpha,\beta)$ has zero loss on $I$, hence is Bayes on
$\Delta_I$. At a prior outside $\Delta_T$ the
Bayes model has $I(\alpha)=\varnothing$ and differs from $M(\alpha,\beta)$
on an open set. At $r\in\Delta_T$ with some $i\in\supp r\setminus I$,
$M(\alpha,\beta)$ accepts an open set of gambles with $g_i<0$, so
$L_i>0$.
\end{proof}

\begin{example}\label{ex:covered}
The following targets are in $\K$.
\begin{enumerate}[label=(\roman*),leftmargin=*]
\item $V$: one anchored world, all others free with $u=\infty$. A
system in the family covers $V$ only if $|U|\le1$
(Proposition~\ref{prop:V}; compare Corollary~\ref{cor:vertices} below).
\item Every face $\Delta_W$, every precise model ($l=u=m$ on
$\supp m$), and every credal set when $n=2$. Take the interval of
$m\in\Delta_{\{1,2\}}$ with $x\le m_1\le y$. If $y<1$, use $l_2=u_2=1$,
$l_1=x/(1-x)$ and $u_1=y/(1-y)$; world $1$ is free when $x=0$. If $y=1$
and $x>0$, use $l_1=u_1=1$, $l_2=0$ and $u_2=(1-x)/x$. If $x=0$ and $y=1$,
the interval is a face. All other coordinates have $l=u=0$.
\item The sets
$\{m:m_3\ge\theta\max(m_1/p_1,m_2/p_2)\}$, $\theta>0$, of Example 5.1 in
\texttt{truth\_\allowbreak directed\_\allowbreak losses.tex}, where
$\theta$ is the constant written $\eta$ there. Here world $3$ is anchored,
and worlds $1$ and $2$ are free with $u_i=p_i/\theta$. Their boundary points have
incomparable supports $\{1,3\}$ and $\{2,3\}$.
\item Single-anchor sets $\{m\in\Delta:m_b\le c_bm_a\ (b\ne a)\}$ with
$c_b\in[0,\infty]$.
\item $\conv(K\cup\Delta_I)$, for a normalized box $K$ whose lower bounds
vanish on $I$: raise those upper bounds to $\infty$.
\end{enumerate}
\end{example}

\begin{proposition}\label{prop:notcovered}
Let $n\ge3$. The following targets are not in $\K$, so no $\Sigma_T$
covers them:
\begin{enumerate}[label=(\alph*),leftmargin=*]
\item every $\varepsilon$-contamination class $(1-\varepsilon)p+\varepsilon\Delta$
with $0<\varepsilon<1$;
\item every segment from an interior $p$ to a vertex $e_j$;
\item for $n=4$ and interior $p$, $\conv(\Delta_{\{1,2,3\}}\cup\{p\})$,
a proper target containing three vertices.
\end{enumerate}
More generally, every facet of a full-dimensional member of $\K$ lies in a
hyperplane $\{m_k=0\}$ or $\{m_b=\rho m_a\}$. So a full-dimensional
polytope with any other facet is not in $\K$, and no $\Sigma_T$ covers it.
\end{proposition}
\begin{proof}
A full-dimensional $K\in\K$ has $P=\Omega$, and by
Lemma~\ref{lem:intrinsic} it is the solution set in the affine hull of
$\Delta$ of finitely many inequalities $c\cdot m\le0$, each with
$c=l_ae_b-u_be_a$ or $c=-e_k$. Every facet is cut out by one of them.
Within $\{\one\cdot m=1\}$, a hyperplane has a homogeneous equation $c\cdot
m=0$ that is unique up to a nonzero factor: if $c'\cdot m=0$ defines the
same hyperplane, then $c'=\gamma c+\tau\one$ and $\tau=0$. If the
inequality is oriented so that $K\subseteq\{c\cdot m\le0\}$, the factor is
positive. Each listed $c$ has at most one negative and at most one
positive entry.

(a) For $p_i>0$, the set $\{m_i=(1-\varepsilon)p_i\}$ is a facet. Its
homogeneous normal $e_i-(1-\varepsilon)p_i\one$ has $n-1\ge2$ negative
entries.

(c) Let $\{a,b,c\}=\{1,2,3\}$. The facet through $e_a,e_b,p$ lies in
$\{p_cm_4=p_4m_c\}$. Since $e_c\in K$ gives $p_c\cdot0-p_4\cdot1<0$, the
oriented inequality is $p_cm_4-p_4m_c\le0$. Its normal $p_ce_4-p_4e_c$ is a
positive multiple of a listed $l_{a'}e_{b'}-u_{b'}e_{a'}$ only if $b'=4$ and
$a'=c$. This forces $c\in L$. Hence $L\supseteq\{1,2,3\}$,
so $u_1<\infty$ and $L\not\subseteq\{1\}$. By Lemma~\ref{lem:intrinsic},
$e_1\notin K$.

(b) Suppose $K(l,u)$ is such a segment. Then $P=\Omega$, $L\ne\varnothing$,
and by Lemma~\ref{lem:intrinsic} either $u_j=\infty$ or $L=\{j\}$. If
$u_j=\infty$, then $l_j=0$, and box points with $a_j=0$ give a point of $K$
with $m_j=0$. If $L=\{j\}$, the box contains every $a$ with $a_j=l_j$ and
$0<a_b<u_b$ for $b\ne j$. Normalization is injective on the hyperplane
$\{a_j=l_j\}$, so $K$ has dimension $n-1\ge2$. Either way $K$ is not
the segment.
\end{proof}

\section{Comparison with the three-vertex obstruction}

The additive obstruction of \emph{Full Coverage by Homothetic Credal
Families v3} states that, under (LA) and (S), every target containing
three vertices is strictly dominated. It assumes
$0<\ell_i(g,0)<\ell_i(g,\text{bad})$ at every world, off null sets. In
$\Sigma_T$ the first inequality fails at every world of $T$ on mixed
problems. The proof of every step that uses it no longer applies, and the
first step is in fact false in $\Sigma_T$:
\begin{itemize}[leftmargin=*,itemsep=2pt]
\item Part (c) of the first lemma in its Appendix C, that a coherent
$p$-Bayes model has its credal set in $\Delta_{\supp p}$, fails at priors in $\Delta_T$, where these are
ties (Proposition~\ref{prop:bayes}(b)). The covering models with
$I(\alpha)\ne\varnothing$ are Bayes on $\Delta_{I(\alpha)}$, yet their
credal sets leave that face.
\item The step that finds ``a third index with positive weight and positive
silence penalty'' needs a strictly scored vertex world.
\item The three-world compatibility lemma needs $\operatorname{sgn}
\psi_i=\operatorname{sgn}g_i$, while $\psi_i=0$ when $g_i>0$ and $i\in T$.
\end{itemize}
The finitely additive version in \emph{Closed Choice Beyond Ordinary
Integrals} assumes a strictly guidance-respecting extension, and fails to
apply for the same reason.

The family covers exactly the three-vertex targets in $\K$, by
Lemma~\ref{lem:intrinsic}: those with at least three indices $k\in P$
such that $u_k=\infty$ or $L\subseteq\{k\}$. Examples in four worlds are the proper face
$\Delta_{\{1,2,3\}}$ and the full-dimensional set
$\{m\in\Delta:m_4\le c\,m_3\}$. It does not cover
$\conv(\Delta_{\{1,2,3\}}\cup\{p\})$ (Proposition~\ref{prop:notcovered}).
The obstruction is respected in a sharper form within the family.

\begin{corollary}\label{cor:vertices}
If $\M_K$ is undominated in $\Sigma_T$ and $K$ contains the vertices $e_k$
for $k\in E$, then $|E\setminus T|\le1$. In particular, a three-vertex
target is covered only by systems that conditionally shun error at two or
more of its vertex worlds.
\end{corollary}
\begin{proof}
If $U=\varnothing$, then $E\setminus T=\varnothing$. Otherwise, by
Theorem~\ref{thm:undominated} and Lemma~\ref{lem:credal}, $\M_K$ agrees
almost everywhere with $\M_{K(l,u)}$ for the pair $(l,u)$ of
Lemma~\ref{lem:credal}, and Lemma~\ref{lem:ident} gives $K=K(l,u)$. There
$L=\supp\beta\ne\varnothing$, and $u_k=\infty$ exactly when
$k\in I(\alpha)\subseteq T$. By Lemma~\ref{lem:intrinsic}, $e_k\in K$
requires $k\in I(\alpha)$ or $\supp\beta=\{k\}$. The second can hold for at
most one $k$.
\end{proof}

\section{Status}

\begin{itemize}[leftmargin=*,itemsep=3pt]
\item \textbf{Proved.} For every $T$, $\lambda$ and $\eta$, $\Sigma_T$ has
finite complete rankings satisfying the weakened (G), (IB), (LA), (C) and
(S). Its undominated models are exactly the $M(\alpha,\beta)$ (only $V$
when $U=\varnothing$), and $V$ is among them if and only if $|U|\le1$. The family covers exactly $\K$.
\item \textbf{Proved, negative.} Bayes-at-the-prior selections fail (S)
unless $T\in\{\varnothing,\Omega\}$ or $n=2$. The variant whose
$T$-worlds credit gambles paying throughout $T$ can be incoherent, and
then violates (S) (Remark~\ref{rem:variant}).
\item \textbf{Not covered by the family.} Targets outside $\K$, including
$\varepsilon$-contamination classes, segments from interior points to
vertices, and $\conv(\Delta_{\{1,2,3\}}\cup\{p\})$. In the parametrization
above, a box coordinate with $l_i>0$ and $u_i=\infty$ would correspond to
$\lambda_i>0$ and $\eta_i=\infty$, an infinite penalty for bad advice.
Remark 5.3 of \texttt{truth\_\allowbreak directed\_\allowbreak losses.tex} shows that such
penalties break (S) in a closely related setting. This is a heuristic, not
an impossibility theorem.
\item \textbf{Open.} Whether an additive system outside this family, with
strict penalties at two of three vertex worlds, can cover a three-vertex
target. Also, whether any system satisfying the axioms, with (G) weakened
only by conditional shunning, covers the targets listed above.
\end{itemize}
\end{document}
